% GENERATED FILE. Edit canonical lessons, then run npm run generate:books.
% canonical-source-hash: fnv1a64:fe60ecdd9f278069
% canonical-lessons: FR-C326-liberer, FR-C326-liberation, FR-C326-resister, FR-C326-resistance, FR-C326-commemorer

\chapter{La libération}
\label{ch:fr-la-liberation}
\hlchaptermodality{326}

\begin{chapteropening}
\textbf{By the end of this chapter:} \emph{I can talk about freedom and remembrance: freeing, liberation, resisting, the resistance and commemorating.}

Complete the last lesson of chapter 326: I can talk about freedom and remembrance: freeing, liberation, resisting, the resistance and commemorating.
\end{chapteropening}

\section[libérer]{libérer — to free, to set free}

\label{lesson:FR-C326-liberer}



\begin{quote}
Before the new one: say the French for the archaeologist, then the French for the dig.
\end{quote}

\subsection*{You'll want to know: libérer}
\textbf{libérer} — \textquotedblleft{}to free, to set free\textquotedblright{}.

\begin{quote}
Les soldats ont libéré la ville.
\end{quote}

\begin{quote}
\emph{The soldiers freed the town.}
\end{quote}

Also of a room: \emph{libérer la chambre} — to check out.

\begin{cousinweb}
\emph{libérer} is a learned borrowing from Latin \emph{līberāre} \textquotedblleft{}to set free\textquotedblright{} (EtymDB), built on \emph{līber} \textquotedblleft{}free\textquotedblright{}.
\end{cousinweb}

\begin{grammarlens}[title={What you've built}]
One more word for telling, explaining and describing at B1.
\end{grammarlens}

\subsection*{Your turn}
\begin{itemize}
\raggedright
  \item \emph{Say it:} \emph{libérer}
  \item \emph{Say it:} \emph{libérer}, once more
  \item \emph{Say it:} a sentence of your own with \emph{libérer}
  \item \emph{Recall:} say the French for the archaeologist, then the French for the dig, then say \emph{libérer} again
\end{itemize}

\begin{tcolorbox}[breakable,colback=teal!4,colframe=teal!35!black,title={Before you move on}]
What does libérer mean? (\textquotedblleft{}To free, to set free\textquotedblright{}.) Say it once more.
\end{tcolorbox}
\section[la libération]{la libération — the liberation, the release}

\label{lesson:FR-C326-liberation}



\begin{quote}
Before the new one: say the French for the dig, then the French for to free.
\end{quote}

\subsection*{You'll want to know: la libération}
\textbf{la libération} — \textquotedblleft{}the liberation, the release\textquotedblright{}.

\begin{quote}
La libération de la ville a été une grande fête.
\end{quote}

\begin{quote}
\emph{The liberation of the town was a great celebration.}
\end{quote}

\emph{La libération d'un otage} — the release of a hostage.

\begin{cousinweb}
\emph{libération} is a learned borrowing from Latin \emph{līberātiō} \textquotedblleft{}a setting free, a release\textquotedblright{} (EtymDB).
\end{cousinweb}

\begin{grammarlens}[title={What you've built}]
One more word for telling, explaining and describing at B1.
\end{grammarlens}

\subsection*{Your turn}
\begin{itemize}
\raggedright
  \item \emph{Say it:} \emph{la libération}
  \item \emph{Say it:} \emph{la libération}, once more
  \item \emph{Say it:} a sentence of your own with \emph{la libération}
  \item \emph{Recall:} say the French for the dig, then the French for to free, then say \emph{la libération} again
\end{itemize}

\begin{tcolorbox}[breakable,colback=teal!4,colframe=teal!35!black,title={Before you move on}]
What does la libération mean? (\textquotedblleft{}The liberation, the release\textquotedblright{}.) Say it once more.
\end{tcolorbox}
\section[résister]{résister — to resist, to hold out}

\label{lesson:FR-C326-resister}



\begin{quote}
Before the new one: say the French for to free, then the French for the liberation.
\end{quote}

\subsection*{You'll want to know: résister}
\textbf{résister} — \textquotedblleft{}to resist, to hold out\textquotedblright{}.

\begin{quote}
Le petit village a résisté pendant des mois.
\end{quote}

\begin{quote}
\emph{The little village held out for months.}
\end{quote}

\emph{Résister à} — to resist: \emph{résister à l'ennemi}.

\begin{cousinweb}
\emph{résister} is a learned borrowing from Latin \emph{resistere} \textquotedblleft{}to stand back, to withstand\textquotedblright{} (EtymDB), built on \emph{sistere}, from \emph{stāre} \textquotedblleft{}to stand\textquotedblright{}.
\end{cousinweb}

\begin{grammarlens}[title={What you've built}]
One more word for telling, explaining and describing at B1.
\end{grammarlens}

\subsection*{Your turn}
\begin{itemize}
\raggedright
  \item \emph{Say it:} \emph{résister}
  \item \emph{Say it:} \emph{résister}, once more
  \item \emph{Say it:} a sentence of your own with \emph{résister}
  \item \emph{Recall:} say the French for to free, then the French for the liberation, then say \emph{résister} again
\end{itemize}

\begin{tcolorbox}[breakable,colback=teal!4,colframe=teal!35!black,title={Before you move on}]
What does résister mean? (\textquotedblleft{}To resist, to hold out\textquotedblright{}.) Say it once more.
\end{tcolorbox}
\section[la résistance]{la résistance — the resistance; stamina}

\label{lesson:FR-C326-resistance}



\begin{quote}
Before the new one: say the French for the liberation, then the French for to resist.
\end{quote}

\subsection*{You'll want to know: la résistance}
\textbf{la résistance} — \textquotedblleft{}the resistance; stamina\textquotedblright{}.

\begin{quote}
Ses grands-parents ont fait partie de la Résistance.
\end{quote}

\begin{quote}
\emph{His grandparents were part of the Resistance.}
\end{quote}

Also stamina: \emph{avoir de la résistance}.

\begin{cousinweb}
\emph{résistance} is built on \emph{résister}, a learned borrowing from Latin \emph{resistere} \textquotedblleft{}to withstand\textquotedblright{} (Lewis \& Short).
\end{cousinweb}

\begin{grammarlens}[title={What you've built}]
One more word for telling, explaining and describing at B1.
\end{grammarlens}

\subsection*{Your turn}
\begin{itemize}
\raggedright
  \item \emph{Say it:} \emph{la résistance}
  \item \emph{Say it:} \emph{la résistance}, once more
  \item \emph{Say it:} a sentence of your own with \emph{la résistance}
  \item \emph{Recall:} say the French for the liberation, then the French for to resist, then say \emph{la résistance} again
\end{itemize}

\begin{tcolorbox}[breakable,colback=teal!4,colframe=teal!35!black,title={Before you move on}]
What does la résistance mean? (\textquotedblleft{}The resistance; stamina\textquotedblright{}.) Say it once more.
\end{tcolorbox}
\section[commémorer]{commémorer — to commemorate}

\label{lesson:FR-C326-commemorer}



\begin{quote}
Before the new one: say the French for to resist, then the French for the resistance.
\end{quote}

\subsection*{You'll want to know: commémorer}
\textbf{commémorer} — \textquotedblleft{}to commemorate\textquotedblright{}.

\begin{quote}
Chaque année, la ville commémore la fin de la guerre.
\end{quote}

\begin{quote}
\emph{Every year, the town commemorates the end of the war.}
\end{quote}

The noun is \emph{la commémoration}.

\begin{cousinweb}
\emph{commémorer} is a learned borrowing from Latin \emph{commemorāre} \textquotedblleft{}to call to mind\textquotedblright{} (EtymDB), from \emph{con-} and \emph{memor} \textquotedblleft{}mindful\textquotedblright{}.
\end{cousinweb}

\begin{grammarlens}[title={What you've built}]
One more word for telling, explaining and describing at B1.
\end{grammarlens}

\subsection*{Your turn}
\begin{itemize}
\raggedright
  \item \emph{Say it:} \emph{commémorer}
  \item \emph{Say it:} \emph{commémorer}, once more
  \item \emph{Say it:} a sentence of your own with \emph{commémorer}
  \item \emph{Recall:} say the French for to resist, then the French for the resistance, then say \emph{commémorer} again
\end{itemize}

\begin{tcolorbox}[breakable,colback=teal!4,colframe=teal!35!black,title={Before you move on}]
What does commémorer mean? (\textquotedblleft{}To commemorate\textquotedblright{}.) Say it once more.
\end{tcolorbox}
